% GENERATED by build_chapters.py; edit docs/report/chapters, not this file

\chapter{System Architecture}

This chapter describes the whole system from above: which parts exist, what each one owns, how they are wired together, and why the main design choices were made. Chapter 5 covers the GitHub side in detail and Chapter 6 covers the n8n workflow.

\section{Components}

\begin{figure}[!htbp]
\centering
\begin{tikzpicture}[every node/.style={font=\sffamily\footnotesize},
    box/.style={rectangle, rounded corners=3pt, draw=black!70, fill=blue!8, minimum width=3.5cm, minimum height=1.1cm, align=center},
    gbox/.style={box, fill=green!10}, rbox/.style={box, fill=orange!12}]
  \node[gbox] (push) at (0,0)     {Developer pushes\\to \texttt{main}};
  \node[box]  (ci)   at (4.5,0)   {GitHub Actions\\\texttt{build-and-test}};
  \node[box]  (note) at (9.0,0)   {\texttt{notify} job\\(webhook call)};
  \node[box]  (n8n)  at (13.5,0)  {n8n workflow\\(26 nodes)};
  \node[gbox] (api)  at (9.0,-2.6) {GitHub API\\logs, files, branch, PR};
  \node[rbox] (model) at (13.5,-2.6) {Qwen3.8-27B on Kaggle\\2$\times$T4, via tunnel};
  \node[box]  (back) at (13.5,-5.2) {\texttt{ci-heal-backend}\\(alternative healer)};
  \draw[flow] (push) -- (ci);
  \draw[flow] (ci) -- node[lbl, above]{code failure} (note);
  \draw[flow] (note) -- node[lbl, above]{POST} (n8n);
  \draw[flow] (n8n.south) -- node[lbl, right]{prompt} (model.north);
  \draw[flow] (model.west) -- node[lbl, below]{edits} (api.east);
  \draw[flow] (n8n.south west) -- node[lbl, above, sloped]{fix PR} (api.north east);
  \draw[flow, dashed] (back.north) -- (model.south);
  \draw[flow, dashed] (back.west) -| (api.south);
\end{tikzpicture}
\caption{Components of the system and the path of one heal. The dashed arrows show the second healer, which uses the same model endpoint and the same GitHub API.}
\label{fig:architecture}
\end{figure}


The system has six parts. Two of them (the CI workflow and the n8n workflow) take part in every heal, one (the model endpoint) is shared by everything, and the rest are an alternative healer and support tools.

\begingroup\small
\begin{longtable}{>{\raggedright\arraybackslash}p{0.223\dimexpr\linewidth-6\tabcolsep\relax}>{\raggedright\arraybackslash}p{0.345\dimexpr\linewidth-6\tabcolsep\relax}>{\raggedright\arraybackslash}p{0.432\dimexpr\linewidth-6\tabcolsep\relax}}
\toprule
\textbf{Component} & \textbf{What it owns} & \textbf{Where it lives} \\
\midrule
\endhead
\bottomrule
\endfoot
\endlastfoot
GitHub Actions workflow & Runs the tests, decides whether a failure is a code bug, counts heal attempts, calls the healer & \texttt{.\allowbreak{}github/\allowbreak{}workflows/\allowbreak{}demo-\allowbreak{}ci.\allowbreak{}yml} \\
n8n workflow & Collects the evidence, asks the model, applies the edits, opens or updates the fix pull request, reports failures & \texttt{n8n/\allowbreak{}self\_\allowbreak{}healing\_\allowbreak{}pipeline.\allowbreak{}json} (26 nodes, generated by \texttt{tools/\allowbreak{}gen\_\allowbreak{}n8n.\allowbreak{}py}) \\
Kaggle endpoint & Serves Qwen3.8-27B over an OpenAI-style chat API on two T4 GPUs, behind a Cloudflare quick tunnel & \texttt{kaggle/\allowbreak{}endpoint/\allowbreak{}deploy\_\allowbreak{}src.\allowbreak{}py} \\
ci-heal-backend & The alternative healer: a FastAPI service that classifies a failed run, then retries it, opens a fix PR or opens an issue & \texttt{ci-\allowbreak{}heal-\allowbreak{}backend/\allowbreak{}} \\
Launcher scripts & Double-click start and stop scripts for the model, n8n and the backend & \texttt{launcher/\allowbreak{}} \\
Evaluation harness & Bug cases, scorers, unit tests and a real-n8n rig that produced the measurements used in this report & \texttt{eval/\allowbreak{}}, \texttt{tools/\allowbreak{}tests/\allowbreak{}}, \texttt{tools/\allowbreak{}n8n\_\allowbreak{}rig/\allowbreak{}}, \texttt{run\_\allowbreak{}tests.\allowbreak{}py} \\
\bottomrule
\caption{Components}\label{tab:04-1}\\
\end{longtable}
\endgroup

(source: README.md, "Repository layout"; docs/\allowbreak{}architecture/\allowbreak{}n8n-workflow.md; docs/\allowbreak{}architecture/\allowbreak{}kaggle-endpoint.md)

The JSON of the n8n workflow is never edited by hand. \texttt{tools/\allowbreak{}gen\_\allowbreak{}n8n.\allowbreak{}py} writes it, and the bodies of its Code nodes are plain files in \texttt{tools/\allowbreak{}nodes/\allowbreak{}}. Running \texttt{python tools/\allowbreak{}gen\_\allowbreak{}n8n.\allowbreak{}py -\allowbreak{}-\allowbreak{}check} regenerates the file in memory and compares it byte for byte with the shipped copy (source: docs/\allowbreak{}architecture/\allowbreak{}n8n-workflow.md, checked against the node list in n8n/\allowbreak{}self\_\allowbreak{}healing\_\allowbreak{}pipeline.json).

The Kaggle endpoint runs as a script kernel. At boot it unpacks a prebuilt \texttt{llama-\allowbreak{}server}, downloads the GGUF model file, tries context sizes from 131072 downward until one boots, starts a Cloudflare quick tunnel and prints the public address (source: docs/\allowbreak{}architecture/\allowbreak{}kaggle-endpoint.md, describing kaggle/\allowbreak{}endpoint/\allowbreak{}deploy\_\allowbreak{}src.py). The script keeps serving for a 1.5 hour window and then ends by itself.

The launcher is a convenience layer. \texttt{1\_\allowbreak{}start\_\allowbreak{}model\_\allowbreak{}on\_\allowbreak{}kaggle.\allowbreak{}bat} checks the kernel status, builds and pushes the kernel, watches the boot log until the tunnel address appears, and then posts that address to the n8n workflow (source: launcher/\allowbreak{}README.md).

The evaluation harness matters for the architecture because the Code nodes can run outside n8n. \texttt{run\_\allowbreak{}tests.\allowbreak{}py} runs the node suites in \texttt{tools/\allowbreak{}tests/\allowbreak{}}, the eval scorer self-test and the generator check, with no GPU and no writes to GitHub. With \texttt{-\allowbreak{}-\allowbreak{}rig} it also runs a real n8n 1.x engine against mock GitHub hosts (source: run\_\allowbreak{}tests.py, docstring).

\section{Data flow of one heal}

\begin{figure}[!htbp]
\centering
\begin{tikzpicture}[every node/.style={font=\sffamily\footnotesize}]
  \node[ext] (dev) at (0,0) {Developer};
  \node[dfdproc, minimum size=3cm] (core) at (6.4,0) {0.0\\Self-Healing\\Pipeline};
  \node[ext] (gh) at (13.2,2.4) {GitHub\\(Actions, API)};
  \node[ext] (llm) at (13.2,-2.4) {Model endpoint\\(Qwen3.8)};
  \draw[flow] ([yshift=0.3cm]dev.east) -- node[lbl, above]{push of broken code} ([yshift=0.3cm]dev.east -| core.west);
  \draw[flow] ([yshift=-0.3cm]dev.east -| core.west) -- node[lbl, below]{fix pull request} ([yshift=-0.3cm]dev.east);
  \draw[flow] (core.50) -- node[lbl, above, pos=0.42]{fix PR, comments} ([yshift=0.35cm]gh.west);
  \draw[flow] ([yshift=-0.35cm]gh.west) -- node[lbl, below, pos=0.55]{events, logs, files} (core.25);
  \draw[flow] (core.-25) -- node[lbl, above, pos=0.42]{prompt} ([yshift=0.35cm]llm.west);
  \draw[flow] ([yshift=-0.35cm]llm.west) -- node[lbl, below, pos=0.55]{edits} (core.-50);
\end{tikzpicture}
\caption{Level 0 data flow diagram: the whole pipeline as one process with its three external actors.}
\label{fig:dfd0}
\end{figure}

\begin{figure}[!htbp]
\centering
\begin{tikzpicture}[every node/.style={font=\sffamily\scriptsize}]
  \node[ext, minimum width=2.4cm] (gh) at (-0.4,0.3) {GitHub\\webhook / API};
  \node[dfdproc, minimum size=2.2cm] (p1) at (4.0,2.2)  {1.0\\Receive and\\authenticate};
  \node[dfdproc, minimum size=2.2cm] (p2) at (8.0,2.2)  {2.0\\Gather\\evidence};
  \node[dfdproc, minimum size=2.2cm] (p3) at (12.0,2.2) {3.0\\Build\\prompt};
  \node[dfdproc, minimum size=2.2cm] (p4) at (12.0,-1.6) {4.0\\Call model\\(fallback)};
  \node[dfdproc, minimum size=2.2cm] (p5) at (8.0,-1.6) {5.0\\Extract\\and check fix};
  \node[dfdproc, minimum size=2.2cm] (p6) at (4.0,-1.6) {6.0\\Open PR or\\comment};
  \node[store] (d1) at (8.0,-4.4) {D1  workflow static data: model URL};
  \draw[flow] (gh) -- node[lbl, above, sloped, pos=0.55]{failure event} (p1);
  \draw[flow] (p1) -- node[lbl, above]{run id} (p2);
  \draw[flow] (p2) -- node[lbl, above]{log, files} (p3);
  \draw[flow] (p3) -- node[lbl, right]{request} (p4);
  \draw[flow] (p4) -- node[lbl, above]{reply} (p5);
  \draw[flow] (p5) -- node[lbl, above]{edits} (p6);
  \draw[flow] (p6) -- node[lbl, below, sloped, pos=0.45]{branch, PR} (gh);
  \draw[flow, dashed] (d1.north) -| node[lbl, pos=0.2, right]{address} (p4.south);
\end{tikzpicture}
\caption{Level 1 data flow diagram: the six sub-processes of the n8n healer and the one data store.}
\label{fig:dfd1}
\end{figure}


The steps below follow a first-time heal from the push to the open pull request.

\begin{enumerate}[leftmargin=*, itemsep=2pt]
\item A developer pushes a commit that changes \texttt{demo-\allowbreak{}app/\allowbreak{}} to \texttt{main}, and the workflow starts (source: .github/\allowbreak{}workflows/\allowbreak{}demo-ci.yml, \texttt{on.\allowbreak{}push}).
\item The \texttt{build-\allowbreak{}and-\allowbreak{}test} job runs \texttt{npm test} and turns the exit code into \texttt{failure\_\allowbreak{}kind}: \texttt{none} for 0, \texttt{infra} for 75 and \texttt{code} for anything else (source: .github/\allowbreak{}workflows/\allowbreak{}demo-ci.yml, step "Run tests").
\item The \texttt{notify-\allowbreak{}n8n-\allowbreak{}on-\allowbreak{}failure} job runs only when \texttt{failure\_\allowbreak{}kind} is \texttt{code} and the event is a push or the branch starts with \texttt{ai-\allowbreak{}fix-\allowbreak{}}. On a push it sets \texttt{attempt=0} and \texttt{action=heal}.
\item The notify job POSTs a JSON body to the n8n webhook. The fields are in the table below.
\item n8n answers GitHub at once (response mode \texttt{onReceived}) and keeps working (source: n8n/\allowbreak{}self\_\allowbreak{}healing\_\allowbreak{}pipeline.json, node "GitHub Failure Webhook").
\item n8n lists the run's jobs, finds the failed one, downloads its log, lists the commit's changed files and picks up to 3 candidate files (source: tools/\allowbreak{}nodes/\allowbreak{}pick\_\allowbreak{}candidate\_\allowbreak{}files.js, \texttt{MAX = 3}, as described in docs/\allowbreak{}architecture/\allowbreak{}n8n-workflow.md).
\item n8n builds one prompt from the failure details, the trimmed log and the source of each candidate file. For a first heal it sets \texttt{enable\_\allowbreak{}thinking} to false.
\item n8n sends the request to the model address it has stored. The reply comes back as a stream of server-sent events.
\item n8n rebuilds the reply, reads the JSON object, applies the SEARCH/\allowbreak{}REPLACE blocks to the file the model named and checks that the result parses.
\item n8n creates the \texttt{ai-\allowbreak{}fix-\allowbreak{}.\allowbreak{}.\allowbreak{}.\allowbreak{}} branch from the failing commit, pushes the patched file with a commit message that ends in \texttt{Heal-\allowbreak{}Attempt:\allowbreak{} 1}, and opens a pull request into the branch that failed (source: tools/\allowbreak{}nodes/\allowbreak{}extract\_\allowbreak{}ai\_\allowbreak{}response.js; n8n/\allowbreak{}self\_\allowbreak{}healing\_\allowbreak{}pipeline.json, node "Open Pull Request").
\item The workflow also triggers on pull requests into \texttt{main}, so the fix PR gets its own CI run. If that run fails with a code failure, the notify job runs again with \texttt{attempt=1} and the loop continues (Chapters 5 and 6).
\end{enumerate}

The webhook body has these fields (source: .github/\allowbreak{}workflows/\allowbreak{}demo-ci.yml, step "Trigger n8n webhook").

\begingroup\small
\begin{longtable}{>{\raggedright\arraybackslash}p{0.207\dimexpr\linewidth-6\tabcolsep\relax}>{\raggedright\arraybackslash}p{0.397\dimexpr\linewidth-6\tabcolsep\relax}>{\raggedright\arraybackslash}p{0.397\dimexpr\linewidth-6\tabcolsep\relax}}
\toprule
\textbf{Field} & \textbf{Value} & \textbf{Used for} \\
\midrule
\endhead
\bottomrule
\endfoot
\endlastfoot
\texttt{run\_\allowbreak{}id} & \texttt{github.\allowbreak{}run\_\allowbreak{}id} & Listing the jobs of the failed run \\
\texttt{repo} & \texttt{owner/\allowbreak{}name} & Every GitHub API URL \\
\texttt{branch} & The pushed branch, or the head branch of the fix PR & Base of the new PR; retry detection by the \texttt{ai-\allowbreak{}fix-\allowbreak{}} prefix \\
\texttt{commit} & The pushed commit on a push, the PR head commit on a pull request & Source of file contents; base of the new branch \\
\texttt{actor} & \texttt{github.\allowbreak{}actor} & Shown in the prompt \\
\texttt{failure\_\allowbreak{}kind} & \texttt{code} in practice & A second check inside n8n \\
\texttt{attempt} & \texttt{0} on a push, the \texttt{Heal-\allowbreak{}Attempt} number on a fix PR & Turns thinking on; numbers the next commit \\
\texttt{pr\_\allowbreak{}number} & The PR number on a pull request event & Where a failure comment is posted \\
\bottomrule
\caption{Data flow of one heal}\label{tab:04-2}\\
\end{longtable}
\endgroup

All values are sent as strings inside the JSON. On a push event the expression for \texttt{pr\_\allowbreak{}number} has nothing to read, so it should arrive empty. That is my inference from the expression, not something we observed. \texttt{tools/\allowbreak{}nodes/\allowbreak{}explain\_\allowbreak{}failure.\allowbreak{}js} treats anything that is not all digits as "no PR".

\section{Trust boundaries and secrets}

Four boundaries are crossed during one heal, and each has its own protection.

\begingroup\small
\begin{longtable}{>{\raggedright\arraybackslash}p{0.295\dimexpr\linewidth-4\tabcolsep\relax}>{\raggedright\arraybackslash}p{0.705\dimexpr\linewidth-4\tabcolsep\relax}}
\toprule
\textbf{Boundary} & \textbf{What protects it} \\
\midrule
\endhead
\bottomrule
\endfoot
\endlastfoot
GitHub runner to n8n webhook & A shared secret in the header \texttt{X-\allowbreak{}Webhook-\allowbreak{}Secret}, checked by n8n's header authentication. A wrong or missing secret is rejected before the workflow starts (source: docs/\allowbreak{}architecture/\allowbreak{}n8n-workflow.md) \\
n8n to the GitHub API & A bearer token held in the n8n credential "GitHub PAT", used by nine request nodes (source: docs/\allowbreak{}guides/\allowbreak{}cloud-n8n.md) \\
n8n to the model & A bearer key held in the n8n credential "Qwen3.8 Endpoint Key", checked by \texttt{llama-\allowbreak{}server} through \texttt{-\allowbreak{}-\allowbreak{}api-\allowbreak{}key} \\
n8n to the log storage host & No credential at all, on purpose (see the design decision on redirects) \\
\bottomrule
\caption{Trust boundaries and secrets}\label{tab:04-3}\\
\end{longtable}
\endgroup

The secrets live in these places. No value is printed here.

\begingroup\small
\begin{longtable}{>{\raggedright\arraybackslash}p{0.218\dimexpr\linewidth-4\tabcolsep\relax}>{\raggedright\arraybackslash}p{0.782\dimexpr\linewidth-4\tabcolsep\relax}}
\toprule
\textbf{Secret} & \textbf{Where it lives} \\
\midrule
\endhead
\bottomrule
\endfoot
\endlastfoot
Webhook shared secret & GitHub repository secret \texttt{N8N\_\allowbreak{}WEBHOOK\_\allowbreak{}SECRET}; the n8n credential "N8N Webhook Secret"; the field \texttt{webhook\_\allowbreak{}secret} in \texttt{\%LOCALAPPDATA\%\textbackslash{}selfheal-\allowbreak{}n8n\textbackslash{}secrets.\allowbreak{}json} (source: docs/\allowbreak{}guides/\allowbreak{}cloud-n8n.md) \\
Webhook address & GitHub repository secret \texttt{N8N\_\allowbreak{}WEBHOOK\_\allowbreak{}URL} (source: .github/\allowbreak{}workflows/\allowbreak{}demo-ci.yml) \\
GitHub PAT & The n8n credential "GitHub PAT". For a local n8n, \texttt{tools/\allowbreak{}golive/\allowbreak{}n8n\_\allowbreak{}import.\allowbreak{}py} takes the token from \texttt{gh auth token}, writes it into n8n's credential store and deletes the plaintext import file. On n8n Cloud it is created by hand in the browser (source: docs/\allowbreak{}guides/\allowbreak{}cloud-n8n.md) \\
n8n encryption key & The field \texttt{encryption\_\allowbreak{}key} in \texttt{secrets.\allowbreak{}json}, used by the local n8n (source: tools/\allowbreak{}golive/\allowbreak{}n8n\_\allowbreak{}import.py) \\
Endpoint API key & A constant \texttt{API\_\allowbreak{}KEY} in \texttt{kaggle/\allowbreak{}endpoint/\allowbreak{}deploy\_\allowbreak{}src.\allowbreak{}py}, and the n8n credential "Qwen3.8 Endpoint Key" (source: docs/\allowbreak{}architecture/\allowbreak{}kaggle-endpoint.md, "API key") \\
Current model address & \texttt{\%LOCALAPPDATA\%\textbackslash{}selfheal-\allowbreak{}n8n\textbackslash{}qwen\_\allowbreak{}url.\allowbreak{}txt} on the PC, and the workflow's static data in n8n (source: tools/\allowbreak{}nodes/\allowbreak{}store\_\allowbreak{}model\_\allowbreak{}url.js) \\
Backend secrets & \texttt{ci-\allowbreak{}heal-\allowbreak{}backend/\allowbreak{}.\allowbreak{}env} (\texttt{GITHUB\_\allowbreak{}TOKEN}, \texttt{GITHUB\_\allowbreak{}WEBHOOK\_\allowbreak{}SECRET}, \texttt{N8N\_\allowbreak{}SHARED\_\allowbreak{}TOKEN} and others). \texttt{.\allowbreak{}gitignore} lists \texttt{*.\allowbreak{}env} and \texttt{secrets.\allowbreak{}json} \\
\bottomrule
\caption{Trust boundaries and secrets}\label{tab:04-4}\\
\end{longtable}
\endgroup

The folder \texttt{\%LOCALAPPDATA\%\textbackslash{}selfheal-\allowbreak{}n8n} is outside the repository on purpose. It also holds the local n8n install and database, the tunnel logs and the process ids (source: tools/\allowbreak{}golive/\allowbreak{}up.py, docstring; tools/\allowbreak{}golive/\allowbreak{}exec\_\allowbreak{}stats.py).

Two weak points should be stated plainly. The endpoint key is a fixed string inside a file that is committed to the repository, so anyone with repository access can read it (source: docs/\allowbreak{}architecture/\allowbreak{}kaggle-endpoint.md). And the quick tunnel makes the model server reachable from the internet, so that key is the only thing between a stranger and the GPU.

The JSON of the workflow holds no secrets. \texttt{tools/\allowbreak{}tests/\allowbreak{}test\_\allowbreak{}workflow\_\allowbreak{}structure.\allowbreak{}py} searches it for token prefixes, and the OpenAI node ships with no credential bound (source: docs/\allowbreak{}architecture/\allowbreak{}n8n-workflow.md, "Credentials and variables").

\section{One healer per repository}

The n8n workflow and ci-heal-backend are two separate healers, and a repository should use only one of them. GitHub tells both about the same failed run: n8n through the notify job, the backend through a \texttt{workflow\_\allowbreak{}run} webhook. With both active, each would open its own fix PR for the same bug and the two would compete (source: README.md, "How the parts connect"; launcher/\allowbreak{}README.md).

They also decide differently. The n8n workflow acts only on failures the CI job marked as \texttt{code}, and then always tries a fix. The backend classifies the failure first and then reruns a flaky job, opens a fix PR or opens an issue, depending on its confidence (source: ci-heal-backend/\allowbreak{}README.md). Their attempt counters are separate too: n8n uses the \texttt{Heal-\allowbreak{}Attempt} trailer, while the backend counts fix decisions in its own database (source: ci-heal-backend/\allowbreak{}README.md, Architecture). Neither counter knows about the other.

The backend has an endpoint \texttt{POST /\allowbreak{}webhooks/\allowbreak{}n8n} for failures forwarded from n8n (source: ci-heal-backend/\allowbreak{}README.md). We found no node in the shipped workflow that calls it, so that route is not part of the n8n healer described in this report.

\section{Design decisions}

\subsection{SEARCH/\allowbreak{}REPLACE edits instead of whole files}

The model answers with a JSON object whose \texttt{edits} field holds SEARCH/\allowbreak{}REPLACE blocks. The first version asked for the whole fixed file, which made the model copy nearly all of its input back out. Output tokens cost time on this hardware (the baseline run measured 12.29 tok/\allowbreak{}s), so three formats were measured on 45 calls (source: docs/\allowbreak{}results/\allowbreak{}DIFF\_\allowbreak{}FORMAT\_\allowbreak{}EVAL.md).

\begingroup\small
\begin{longtable}{>{\raggedright\arraybackslash}p{0.358\dimexpr\linewidth-8\tabcolsep\relax}>{\raggedright\arraybackslash}p{0.296\dimexpr\linewidth-8\tabcolsep\relax}>{\raggedright\arraybackslash}p{0.123\dimexpr\linewidth-8\tabcolsep\relax}>{\raggedright\arraybackslash}p{0.222\dimexpr\linewidth-8\tabcolsep\relax}}
\toprule
\textbf{Format} & \textbf{Median completion tokens} & \textbf{Wall time} & \textbf{Applied and parsed} \\
\midrule
\endhead
\bottomrule
\endfoot
\endlastfoot
Whole file & 655 & 56.4 s & 15/\allowbreak{}15 \\
Unified diff & 233 & 23.3 s & 12/\allowbreak{}15 \\
SEARCH/\allowbreak{}REPLACE & 171 & 17.0 s & 15/\allowbreak{}15 \\
\bottomrule
\caption{SEARCH/\allowbreak{}REPLACE edits instead of whole files}\label{tab:04-5}\\
\end{longtable}
\endgroup

(source: docs/\allowbreak{}results/\allowbreak{}DIFF\_\allowbreak{}FORMAT\_\allowbreak{}EVAL.md, "Results")

The unified diff lost on correctness as well as speed. In one scenario the model wrote a patch whose header claimed an added line while the only added line was empty, and a tolerant \texttt{-\allowbreak{}-\allowbreak{}recount} retry applied it as a no-op that still did not parse (source: same file). SEARCH/\allowbreak{}REPLACE avoids this because every block must match the current file exactly once, or the whole fix is rejected. The rule is enforced in \texttt{tools/\allowbreak{}nodes/\allowbreak{}extract\_\allowbreak{}ai\_\allowbreak{}response.\allowbreak{}js}. In the first live round the median model phase fell from about 57 s to 11.4 s (source: docs/\allowbreak{}results/\allowbreak{}GOLIVE\_\allowbreak{}RESULTS.md).

\subsection{Streamed model calls}

The model sits behind a Cloudflare quick tunnel, which cuts a request when the origin stays silent for 120 seconds. A thinking retry once died that way, with HTTP 524 at 125.9 s. A proof run sent the same long request through the same tunnel twice: without streaming it failed with 524 at 125.9 s, and with streaming it finished with HTTP 200 after 166.5 s (source: docs/\allowbreak{}results/\allowbreak{}GOLIVE\_\allowbreak{}RESULTS.md, Round 5). So the request sets \texttt{stream:\allowbreak{} true}, and a Code node called "Assemble Model Reply" rebuilds the ordinary completion from the event stream. A stream that ends with neither \texttt{[DONE]} nor a finish reason counts as cut off and as a failed call (source: tools/\allowbreak{}nodes/\allowbreak{}assemble\_\allowbreak{}model\_\allowbreak{}reply.js, as described in docs/\allowbreak{}architecture/\allowbreak{}n8n-workflow.md).

\subsection{Thinking off first, on for retries}

With thinking off, the first pass healed 15/\allowbreak{}15 control cases and 12/\allowbreak{}18 hard cases, with a median wall time of 9.1 s. With thinking on and a 1024-token budget it healed 18/\allowbreak{}18 hard cases, including two that no earlier configuration had fixed, but the median rose to 36.8 s (source: docs/\allowbreak{}results/\allowbreak{}SHIPPED\_\allowbreak{}QUALITY\_\allowbreak{}EVAL.md, arms B and C). Always thinking would add about 28 s to every heal, so the workflow turns it on only after CI has rejected the fast fix. \texttt{Build AI Prompt} sets \texttt{enable\_\allowbreak{}thinking} to the value of \texttt{isRetry} (source: tools/\allowbreak{}nodes/\allowbreak{}build\_\allowbreak{}ai\_\allowbreak{}prompt.js). The server itself defaults to thinking off, because the backend can only switch it on per request (source: docs/\allowbreak{}architecture/\allowbreak{}kaggle-endpoint.md, "Why thinking is off by default").

The price is that a hard bug needs one more CI cycle. In the live retry rounds, \texttt{multi\_\allowbreak{}hunk} turned green on attempt 2 in 143 to 161 s, and \texttt{callback\_\allowbreak{}arity}, which failed in earlier rounds, turned green on attempt 3 in 249 s (source: docs/\allowbreak{}results/\allowbreak{}GOLIVE\_\allowbreak{}RESULTS.md, Rounds 3 and 5).

\subsection{The failure\_\allowbreak{}kind gate}

An infrastructure problem, such as a port that is already in use, has nothing to fix in the source. If the model were asked anyway, it would invent a patch. The demo tests therefore exit with 75 for environment problems and with 1 for code failures, and the workflow maps these to \texttt{infra} and \texttt{code} (source: demo-app/\allowbreak{}test.js, header comment). We checked on this machine that \texttt{npm test} exits with 75 when port 5000 is already taken. The gate exists twice: the notify job runs only for \texttt{code}, and \texttt{Build AI Prompt} throws \texttt{Refusing to run} for any other non-empty \texttt{failure\_\allowbreak{}kind}, because the webhook is a public address that anyone could post to (source: tools/\allowbreak{}nodes/\allowbreak{}build\_\allowbreak{}ai\_\allowbreak{}prompt.js).

\subsection{Log redirect handling}

GitHub answers a job-log request with a 302 redirect to a signed address on another host. The first design let one HTTP node follow the redirect, assuming the client drops the \texttt{Authorization} header when the host changes. The real-n8n rig showed that its mock storage host received the bearer token. The fix uses two nodes. "Fetch Failed Job Log" does not follow redirects and returns the \texttt{Location}, and "Download Job Log" fetches that address with no credential. The rig now asserts that the token reached the API host and did not reach the storage host (source: docs/\allowbreak{}results/\allowbreak{}GOLIVE\_\allowbreak{}RESULTS.md, Round 2, "Two defects the verification caught"; n8n/\allowbreak{}self\_\allowbreak{}healing\_\allowbreak{}pipeline.json).

\subsection{Model address stored through a webhook}

The quick tunnel gets a new address on every Kaggle boot, so the address has to be delivered again each time. Earlier the update went through the n8n API (source: git commit 4833c44, whose message says it switched from the API to a webhook). On an n8n Cloud trial the REST API is switched off (source: docs/\allowbreak{}guides/\allowbreak{}cloud-n8n.md, "Trial limits"), so the workflow got a second webhook, "Set Model URL", protected by the same shared secret, and the launcher posts the new address to it. The node "Store Model URL" accepts only an address of the form \texttt{https:\allowbreak{}/\allowbreak{}/\allowbreak{}\textless{}name\textgreater{}.\allowbreak{}trycloudflare.\allowbreak{}com} and saves it in the workflow's static data, so a mistyped or hostile value is refused (source: tools/\allowbreak{}nodes/\allowbreak{}store\_\allowbreak{}model\_\allowbreak{}url.js). "Get Model URL" reads the stored address before each model call and falls back to the variable \texttt{QWEN\_\allowbreak{}BASE\_\allowbreak{}URL} (source: tools/\allowbreak{}nodes/\allowbreak{}get\_\allowbreak{}model\_\allowbreak{}url.js).

\subsection{The attempt counter lives in the commit}

The cap on retries has to survive a re-run of the workflow and a restart of n8n. For that reason the count is a \texttt{Heal-\allowbreak{}Attempt:\allowbreak{} N} line in the message of every commit the healer pushes, and the notify job reads it from the head commit of the fix PR (source: .github/\allowbreak{}workflows/\allowbreak{}demo-ci.yml, comment on \texttt{MAX\_\allowbreak{}ATTEMPTS}). The cap is enforced in the notify job only. The n8n workflow does not check it (source: docs/\allowbreak{}architecture/\allowbreak{}n8n-workflow.md, "The retry loop").
